% Implementation Section Template for ACM TOSEM
% Target: Application/System design paper with empirical evaluation

\section{Implementation}
\label{sec:implementation}

% Provide a high-level overview of the implementation stack and major components.
% Guidance: Describe the technology stack (e.g., Python, LaTeX, specific libraries).
% Mention the versioning and environment management (e.g., uv, virtualenv).

\subsection{Technology Stack and Environment}
\label{subsec:tech_stack}

% Placeholder for tech stack details.
% Example: The system is implemented in [Language] using [Framework/Libraries].
% Environment management is handled via [Tool].

\subsection{Core Components and Service Architecture}
\label{subsec:core_components}

% Describe the primary modules or services.
% Use figure environments for architecture diagrams if applicable.

\begin{figure}[h]
	\centering
	% \includegraphics[width=\linewidth]{figures/architecture_diagram}
	\caption{Overview of the system architecture and component interactions.}
	\label{fig:arch_overview}
\end{figure}

% Subsections for specific complex components
\subsection{Data Persistence and Normalization}
\label{subsec:persistence}

% Guidance: Discuss how data is stored (YAML, BibTeX, etc.) and normalized.
% Mention deduplication or integrity logic if relevant.

\subsection{External API Integration}
\label{subsec:api_integration}

% Guidance: Document how the system interacts with external services (e.g., arXiv, OpenAlex).
% Mention rate limiting, error handling, and provider abstractions.

\subsection{Deployment and Scalability}
\label{subsec:deployment}

% Guidance: Describe how the system is deployed (e.g., local MCP server, containerized).
% Discuss scalability considerations or performance optimizations.

% Reference placeholder example:
% As noted in~\citep{ref_key}, the integration of...
